\documentclass[11pt,a4paper]{article}

% --- Core packages ---
\usepackage[utf8]{inputenc}
\usepackage[T1]{fontenc}
\usepackage{lmodern}
\usepackage[a4paper,top=1in,bottom=1in,left=1in,right=1in]{geometry}
\usepackage{amsmath,amssymb,amsfonts}
\usepackage{setspace}
\setstretch{1.15}

% --- Lists ---
\usepackage{enumitem}

% --- Tables ---
\usepackage{booktabs}
\usepackage{longtable}
\usepackage{array}

% --- Symbols (check / warn / info) ---
\usepackage{pifont}
\usepackage{xcolor}
\definecolor{okgreen}{rgb}{0.00,0.55,0.00}
\definecolor{warnorange}{rgb}{0.85,0.45,0.00}
\definecolor{infoblue}{rgb}{0.08,0.30,0.60}
\newcommand{\cmark}{\textcolor{okgreen}{\ding{51}}}
\newcommand{\xmark}{\textcolor{warnorange}{\ding{55}}}
\newcommand{\warn}{\textcolor{warnorange}{\ding{115}}}
\newcommand{\info}{\textcolor{infoblue}{\ding{117}}}

% --- Code listings ---
\usepackage{listings}
\definecolor{codebg}{rgb}{0.97,0.97,0.97}
\definecolor{codekw}{rgb}{0.10,0.25,0.70}
\definecolor{codestr}{rgb}{0.55,0.10,0.10}
\definecolor{codecmt}{rgb}{0.25,0.55,0.25}
\lstdefinestyle{py}{
  language=Python,
  backgroundcolor=\color{codebg},
  basicstyle=\ttfamily\small,
  keywordstyle=\color{codekw}\bfseries,
  stringstyle=\color{codestr},
  commentstyle=\color{codecmt}\itshape,
  showstringspaces=false,
  breaklines=true,
  frame=single,
  rulecolor=\color{gray!50},
  tabsize=2,
  columns=fullflexible,
}
\lstset{style=py}

% --- Hyperlinks ---
\usepackage[colorlinks=true,linkcolor=infoblue,citecolor=infoblue,urlcolor=infoblue,
  pdftitle={DAMPS-MMHCL Speedup Completion Guide}]{hyperref}

\title{Speedup Completion Guide for DAMPS-MMHCL\\[0.4em]
\large Closing the Remaining Gaps to 100\% Compliance with the
\textit{DAMPS\_MMHCL\_Training\_Speedup\_Guide\_EN}}
\author{}
\date{}

\begin{document}
\maketitle

\section*{Overall Verdict}
The current implementation does \textbf{not} yet reach 100\% compliance. Based on a
direct inspection of the GitHub repository and the attached Jupyter notebook, the
runtime-level compliance with the \textit{DAMPS\_MMHCL\_Training\_Speedup\_Guide\_EN}
sits at approximately \textbf{85--90\%}.

\section{What Has Been Implemented Correctly}
The following speedup items from the guide are already wired in correctly and are
visible in the latest commit of the repository:

\begin{itemize}[leftmargin=1.4em,itemsep=2pt]
    \item \cmark\ \textbf{cuFFT plan-cache lockdown} \\
          \texttt{torch.backends.cuda.cufft\_plan\_cache[idx].max\_size = 1} is set
          at startup, preventing plan-cache memory leaks.
    \item \cmark\ \textbf{Orthonormal rFFT/irFFT} \\
          \texttt{torch.fft.rfft} and \texttt{torch.fft.irfft} are invoked with
          \texttt{norm="ortho"}, exactly as recommended for the $d{=}64$ embeddings.
    \item \cmark\ \textbf{Dual-path KNN with FAISS-GPU fallback} \\
          \texttt{DualPathKNN} is parameterised with both
          \texttt{faiss\_threshold} and \texttt{faiss\_use\_gpu} arguments inside
          \texttt{train.py}.
    \item \cmark\ \textbf{\texttt{torch.compile} on the DAMPS submodule} \\
          The forward path of the DAMPS module is compiled using
          \texttt{mode="reduce-overhead"} and \texttt{dynamic=True}, while the
          K-NN rebuild step is correctly left uncompiled (matching the caveat
          in the speedup guide).
    \item \cmark\ \textbf{bfloat16 mixed precision via \texttt{torch.autocast}} \\
          The training loop uses
          \texttt{torch.amp.autocast(device\_type=..., dtype=torch.bfloat16, enabled=use\_amp)}
          without a \texttt{GradScaler}, which is the correct pattern for bf16.
    \item \cmark\ \textbf{Optuna TPE Bayesian HPO} \\
          \texttt{scripts/run\_optuna\_hpo.py} is present and ready to be invoked.
    \item \cmark\ \textbf{wandb Sweeps integration} \\
          Both \texttt{scripts/run\_wandb\_sweep.py} and
          \texttt{scripts/wandb\_sweep.yaml} exist in the repository.
    \item \cmark\ \textbf{Paired $t$-test via \texttt{scipy.stats.ttest\_rel}} \\
          Implemented in \texttt{scripts/paired\_ttest.py}.
    \item \cmark\ \textbf{Permutation-FFT aggregator} \\
          \texttt{scripts/\_aggregate\_permutation\_fft.py} is in place, paired
          with notebook cell~5.5.
\end{itemize}

\section{Three Blocking Issues Preventing Full Compliance}

\subsection{The notebook does not enable \texttt{torch.compile}}
\texttt{train.py} already exposes the command-line flags
\texttt{-{}-torch\_compile\_mode} and \texttt{-{}-torch\_compile\_dynamic}, but the
notebook does not pass them into the subprocess command list. As a result, the
25--35\% forward-pass speedup from \texttt{torch.compile} is \textbf{not actually
activated at runtime}, even though the supporting code is fully present.

\textbf{Fix.} Append the following two key/value pairs to the \texttt{cmd} list
that launches \texttt{train.py} from the notebook:

\begin{lstlisting}
"--torch_compile_mode", "reduce-overhead",
"--torch_compile_dynamic", "1",
\end{lstlisting}

\subsection{The notebook does not pass \texttt{-{}-faiss\_use\_gpu 1}}
The notebook currently passes only \texttt{-{}-faiss\_threshold 60000}, but never
\texttt{-{}-faiss\_use\_gpu 1}. For the Amazon Clothing dataset
($N{=}23{,}033 < 60\text{K}$) this has no practical impact, because the chunked
PyTorch path is selected anyway. However, when scaling up to a larger dataset,
the FAISS-GPU fallback will silently remain disabled.

\textbf{Fix.} Add the GPU FAISS flag to the notebook \texttt{cmd} list as well:

\begin{lstlisting}
"--faiss_use_gpu", "1",
\end{lstlisting}

\subsection{\texttt{n\_runs = 3} instead of \texttt{10}}
The notebook is currently configured with \texttt{n\_runs = 3}, which violates the
statistical-reliability requirement shared by both the architecture spec
(\S4) and the Speedup Guide (\S8): final results must be reported across
\textbf{10 seeds}, with confidence intervals and paired $t$-tests.

\textbf{Fix.} Update the notebook constant to:

\begin{lstlisting}
n_runs: int = 10
\end{lstlisting}

\section{Optional Items (Safe to Skip)}

\begin{itemize}[leftmargin=1.4em,itemsep=2pt]
    \item \info\ \textbf{\texttt{torch\_geometric.nn.HypergraphConv} (\S3)} ---
          The repository still uses a custom HGCN implementation. Adopting the
          PyG operator could unlock an additional 25--35\% forward-pass speedup,
          but skipping it does \emph{not} violate the architecture spec.
    \item \info\ \textbf{\texttt{einops} (\S10)} ---
          A pure readability improvement; it has no measurable effect on runtime
          and can safely be omitted.
    \item \info\ \textbf{\texttt{scipy.stats.vonmises} (\S5)} ---
          Not used. The current GPU-vectorised implementation
          $\hat\theta_c=\operatorname{atan2}\!\bigl(\sum\sin R_i,\sum\cos R_i\bigr)$
          is in fact \emph{faster} than the CPU off-loaded path suggested by the
          guide, because it avoids CPU--GPU synchronisation. This deviation is a
          legitimate, beneficial improvement.
\end{itemize}

\section{Compliance Summary}

\begin{longtable}{|p{8.4cm}|c|}
\hline
\textbf{Speedup Guide Item} & \textbf{Status}\\
\hline
\endhead
\hline
\textbf{Speedup Guide Item} & \textbf{Status}\\
\hline
\endfoot
\hline
\endlastfoot
cuFFT plan-cache lock                                      & \cmark \\
\hline
\texttt{torch.fft.rfft} / \texttt{irfft} (\texttt{norm="ortho"}) & \cmark \\
\hline
Dual-path KNN + FAISS-GPU (code level)                     & \cmark \\
\hline
Dual-path KNN + FAISS-GPU (notebook flag)                  & \warn \\
\hline
\texttt{torch.compile} on DAMPS / HGCN (code level)        & \cmark \\
\hline
\texttt{torch.compile} (notebook flag)                     & \warn \\
\hline
PyG \texttt{HypergraphConv}                                & \xmark\ (optional) \\
\hline
von Mises \texttt{atan2}-based MLE                         & \cmark\ (better than guide) \\
\hline
Optuna TPE HPO                                             & \cmark \\
\hline
wandb Sweeps                                               & \cmark \\
\hline
\texttt{scipy.stats.ttest\_rel}                            & \cmark \\
\hline
bfloat16 \texttt{torch.autocast} (no \texttt{GradScaler})  & \cmark \\
\hline
\texttt{einops}                                            & \xmark\ (optional) \\
\hline
Permutation-FFT ablation                                   & \cmark \\
\hline
10 seeds + CI + paired $t$-test                            & \warn \\
\hline
\end{longtable}

\section{Conclusion}
Once the three blocking issues above are patched --- and most importantly,
once \texttt{torch.compile} is actually activated from the notebook --- the
implementation will reach \textbf{full runtime-level compliance} with the
\textit{DAMPS\_MMHCL\_Training\_Speedup\_Guide\_EN}.

\end{document}
